\section{Shared evidence and scientific validation in a causal assistant}
\label{exec23:admit-scope}
A bank analyst may ask whether a fee waiver prevents delinquency, while several
agents retrieve policy manuals, propose confounders and restate the same
customer-service memorandum. Repetition can make one unsupported assumption
appear to have several independent witnesses. A second failure occurs when
a fluent derivation is accepted because a validator assigns it a low risk
score. Neither successful storage nor a reassuring score establishes that the
waiver is identifiable from the available data. The useful product therefore
needs distinct contracts for admitting evidence, proposing a graph, validating
a mathematical argument and releasing a decision. Two recent papers expose
both useful mechanisms and the assumptions that such contracts must retain.

\subsection{What an admission benchmark actually measures}
\label{exec23:admit-memory}
\citet{exec23_epistemic_admission} distinguish a static admission benchmark
from a running shared-memory environment. The static construction imports
public verdicts, adds authored agent and consumer settings, and scores a fixed
gold action. Its 1,200 episodes are divided into 600 training, 200 validation
and 400 test episodes. Agreement with the admission specification is the
target; it is not an empirical measure of emergent propagation. The live
construction uses six agents for six rounds and 100 configurations per model
family, with additional diagnostic and truth-by-copy sets reported separately.
Every live scenario is authored fiction. This makes the input lineage and
scenario truth available to the evaluator, but does not estimate their
frequency in bank operations.

Only promotion writes to shared memory. Later candidates depend on what
earlier policies admitted, so a false-admission fraction is conditional on a
policy-specific stream. If $N_\pi^F$ false candidates arrive under policy
$\pi$ and $A_\pi^F$ are admitted, the ratio
\begin{equation}
 f_\pi=A_\pi^F/N_\pi^F
 \label{exec23:admit-stream}
\end{equation}
is not a comparison on common candidate inputs. For example, ten admissions
from 100 false candidates give $0.10$, whereas five from ten give $0.50$. The
first policy has the smaller conditional rate and the larger number of
admitted false claims. These are stipulated counts, not the paper's results.
A prototype should therefore use both frozen-candidate replay, which measures
a gate on matched inputs, and a closed-loop experiment, which measures the
whole policy-induced stream. It must not combine their denominators.

The paper's propagation damage counts consumer reads while a false belief is
retrievable. This is simulated exposure, not monetary damage. Assertion
adoption is a separate consumer-output measurement. The reported
$0.97$--$0.99$ adoption of an uncontested false belief is conditional on a
memory-only consumer and a selected pool of 553 probes. It is informative
about that constrained interface, not a bank-wide error probability.
Source-type gating performs well in several authored scenarios but admits
false restatements when the source type supplies the authority signal it
trusts. The lineage oracle rejects copied support yet can admit independently
restated falsehoods. Provenance and truth are different properties
\citep{exec23_epistemic_admission}.

An elementary model makes the distinction precise. Suppose each of three
independent witnesses is correct with probability $p=4/5$. A majority is
correct with probability $3p^2(1-p)+p^3=112/125$. If all three repeat one
root witness, the same majority is correct with probability $4/5$. Copy
detection prevents an unjustified increase in confidence; it cannot improve
the root witness's truth probability by itself. These independent rational
calculations are project examples, not a model fitted to the benchmark
\citep{exec23_admission_conformal_reference}.

\subsection{Truth, lineage and usefulness require separate tests}
\label{exec23:admit-validation}
Input-feed lineage is fixed by the live scenario. For generated claims,
retrieval is logged but the derivation link remains agent-declared and can
be incomplete. A citation may mention a false proposition without endorsing
it, while an uncited sentence may assert the same falsehood. The benchmark
therefore separates assertion grading from provenance tracing. Its assertion
grader receives gold propositions; a deployable admission policy does not.
The distinction must survive any reuse of its evaluation design.

The paper also separates blind human labels from labels revised after seeing
the grader. On the assertion sample, the two blind annotators agree more
closely with one another than either does with the model grader. The consumer
matcher misses rephrased values, and its validation sample is balanced by
matcher category rather than drawn to estimate deployment prevalence.
Reported false assertions and adoption therefore depend on the measurement
procedure. The release should retain raw sentences, blind adjudication,
normalization rules and alternative labels instead of treating one grader as
truth. No evaluated admission policy retracts beliefs although the store
offers retraction actions; correction remains untested, not technologically
impossible. The Mem0 and A-MemGuard comparisons concern particular components
and configurations. They do not license an unrestricted conclusion about
every version or deployment of either system
\citep{exec23_epistemic_admission}.

For the bank prototype, each proposed graph edge should retain the business
assertion, original source, observation date, copying lineage, conflicting
evidence and status of human review. An edge such as ``delinquency causes a
servicing call'' must not become ``a servicing call causes delinquency'' merely
because several agents reproduce the same text. A claim-dependency graph
records support between sentences; an intervention graph records assumptions
about mechanisms. The former can organize evidence for the latter, but cannot
identify an intervention by its own topology.

Usefulness must include abstention. Let $p_C$ and $p_W$ denote probabilities
of a correct and a wrong released answer per evaluated request, let an
abstention have zero stipulated gross value, and let $c$ be mean cost per
request. Under assumed benefit $B$ and loss $L$, the expected net is
\begin{equation}
 V=Bp_C-Lp_W-c.
 \label{exec23:admit-value}
\end{equation}
Take $B=10$, $L=100$ and $c=1$ dollars. A policy with
$(p_C,p_W)=(0.60,0.05)$ has net zero dollars; one with
$(0.30,0.01)$ has net one dollar despite answering correctly less often.
These inputs illustrate a decision rule and are not observed bank values.
Compare admission methods at declared coverage, error costs and workload,
then vary those assumptions. Agreement rates or token counts alone cannot
select the better bank policy~\citep{exec23_admission_conformal_reference}.

\subsection{What graph-structured conformal validation can promise}
\label{exec23:admit-conformal}
\citet{exec23_scientific_conformal} propose progressive scientific validation
with branching and an approximate deducibility graph. Their factuality
definition requires both correct claims and justification from prior context
and scientific ground truth. It does not require a relevant answer. Their
approximate graph has semantic ancestor-connectivity and consistency
properties. Acyclicity, reproducible template matching or a plausible
linguistic causal connective does not alone establish these properties.

There is a sound conditional idea in the process-level argument. Let
$Z_1,\ldots,Z_{n+1}$ be exchangeable complete examples and let $g$ be the same
fixed measurable construction for every example. Then
$(Z_i,g(Z_i))_{i=1}^{n+1}$ is exchangeable because applying the same coordinate
map commutes with permutation. This argument permits rich dependence within
one example. It neither assumes that the claims inside a response are
independent nor manufactures exchangeability across examples. Training a
validator, choosing a threshold on the test cases, changing the generation
policy after calibration, or selecting convenient branches requires a
separate argument. Randomized construction is possible if its randomness is
incorporated with the appropriate symmetric law; empirical repeatability
below one does not verify the paper's deterministic special case.

For a conventional score-coverage construction, a finite calibration sample
uses rank $k=\lceil(n+1)(1-\alpha)\rceil$. If $k=n+1$, an infinity sentinel
may be needed; silently choosing the largest finite calibration score changes
the procedure. With $n=5$ and $\alpha=0.1$, the rank is six. Most importantly,
coverage of a score or prediction set must be linked to the desired event
that the released reasoning is valid. That link requires trustworthy
calibration labels, the correct threshold orientation, a nested filtering
family and the stated semantic conditions. The paper's coherent-factuality
proof invokes those properties in addition to exchangeability. A fallible
validator accepting a false statement with score zero is enough to show why
``accepted prefix'' cannot be substituted for ``true prefix'' in a cascade
prevention argument~\citep{exec23_scientific_conformal,exec23_admission_conformal_reference}.

A marginal guarantee is also different from simultaneous reliability. Five
independent answers, each valid with probability $0.9$, have probability
$1-0.9^5=0.40951$ of at least one invalid answer. Adaptive repeated queries
need a specified selection and error-control contract. This calculation
does not refute a marginal conformal result; it identifies the additional
product question~\citep{exec23_admission_conformal_reference}.

\subsection{Ambiguities to repair before certification}
\label{exec23:admit-repairs}
Several passages prevent direct adoption of the paper as a certificate. Its
main pipeline accepts low risk scores, while its branching prose selects
the highest score above threshold. Its minimum-domain-quantile rule must be
interpreted against an explicit event. For ordinary coverage
$\Pr(S\leq q)$, decreasing $q$ is not conservative: two deterministic
domains with scores and valid own-domain quantiles zero and one each have
own-domain coverage one, but using the minimum quantile zero gives coverage
one half under an equal mixture. In contrast, removing more claims can
preserve an all-retained-claims-valid property when the candidate family is
properly nested and validity hereditary. The counterexample concerns the
ordinary score event, not that distinct filtering theorem. The implementation
must state which event it guarantees and prove the relevant nesting rather
than transport the word ``conservative'' between them
\citep{exec23_scientific_conformal,exec23_admission_conformal_reference}.

The displayed progressive-generation pseudocode increments its attempt
counter only when it is below the maximum. If the last attempt again fails,
the loop condition remains true and no exit or counter increment is shown.
An independent bounded trace with maximum two gives attempt states
$1,2,2,2,2$. This is a finding about the printed algorithm, not evidence that
an uninspected implementation hangs. A prototype should use a bounded
attempt loop, distinguish an incomplete retained prefix from a completed
answer, and return an explicit abstention when its budget is exhausted.
The paper's cycle-to-linear-graph fallback likewise needs semantic testing;
turning a graph into a DAG does not establish the claimed deduction relation.

The reported empirical ranking also needs reconciliation. Table 1 gives
MATH accuracy $48.9$ for the descendant-weighted method and $71.2$ for one
listed comparator, despite its universal-superiority caption. Table 2 changes
the full-system MATH and ScienceQA accuracies, and a partial configuration
reports PhyX accuracy $53.1$ and validity $92.1$, compared with $50.1$ and
$91.7$ for the full system. Main-text and appendix model roles, overall
validity, and some domain figures also differ. These observations justify
requesting exact run manifests, denominators, paired outcomes and frozen
calibration sets before ranking methods. They do not establish misconduct,
nor do they erase the useful design ideas. The local review reproduces
arithmetic and control-flow witnesses only; it does not reproduce the
reported model experiments~\citep{exec23_scientific_conformal,exec23_admission_conformal_reference}.

\subsection{A bounded public-evidence prototype route}
\label{exec23:admit-route}
The proposal can proceed without costly bank experiments. Construct public
or authored business scenarios with declared truth and copying lineage; keep
scoring-only labels out of agent inputs. Freeze separate tasks for translating
a business request, proposing competing intervention graphs, identifying a
query, checking an argument and deciding whether to answer. Replay admission
gates on common candidates, then test the entire interacting pipeline with
independent scenario units. Include paraphrase, conflicting authoritative
text, missing lineage, incorrect graph orientation, unsupported deduction,
cyclic dependency output and exhausted validation budgets.

Fit scoring components on training examples; choose the complete generation,
branching and validation procedure before calibration; evaluate on untouched
examples under the declared sampling law. Record refusals, incomplete answers,
truthful corrections, coverage, resource use and conditional errors separately.
Inspect authentic software, licenses and run artifacts before choosing direct
adoption or a port. The memory benchmark is presently a design reference and
the scientific-validation paper a hypothesis and cautionary comparison.
Neither supplies a ready certificate for causal discovery, bank identification
or universal scientific truth. This staged use of public evidence can establish
whether a prototype merits a later bank study while preserving the scientific
questions that the eventual study must answer.
